\documentclass[11pt,a4paper,landscape]{article}
\usepackage[margin=1.2cm]{geometry}
\usepackage{fontspec}
\setmainfont{Times New Roman}
\usepackage{xcolor}
\usepackage{tikz}
\usetikzlibrary{arrows.meta,positioning,fit,calc,backgrounds,shapes.multipart}
\usepackage{amsmath}
\usepackage{amssymb}
\usepackage{array}
\usepackage{booktabs}
\usepackage{graphicx}
\pagestyle{empty}
\setlength{\parindent}{0pt}

\definecolor{inputblue}{HTML}{DCEBFF}
\definecolor{fuzzygreen}{HTML}{DDF5E3}
\definecolor{seqorange}{HTML}{FDE7C8}
\definecolor{mixpurple}{HTML}{E9DDF8}
\definecolor{outred}{HTML}{FADBDD}
\definecolor{notegray}{HTML}{F3F4F6}
\definecolor{linegray}{HTML}{4B5563}
\definecolor{titleblue}{HTML}{1F4E79}

\tikzset{
  >={Latex[length=3mm,width=2mm]},
  line/.style={draw=linegray, -Latex, line width=0.9pt},
  dashedline/.style={draw=linegray, -Latex, line width=0.8pt, dashed},
  box/.style={
    draw=linegray,
    rounded corners=2mm,
    align=center,
    minimum height=11mm,
    minimum width=28mm,
    font=\small,
    inner sep=4pt
  },
  smallbox/.style={
    draw=linegray,
    rounded corners=1.5mm,
    align=center,
    minimum height=9mm,
    minimum width=24mm,
    font=\footnotesize,
    inner sep=3pt
  },
  groupbox/.style={
    draw=linegray,
    rounded corners=2.5mm,
    inner sep=6pt,
    fill=white
  },
  note/.style={
    draw=linegray,
    rounded corners=2mm,
    fill=notegray,
    align=left,
    font=\footnotesize,
    inner sep=5pt,
    text width=9.2cm
  }
}

\begin{document}

\begin{center}
{\LARGE\bfseries\color{titleblue} Sơ đồ tổng quan mô hình \texttt{run\_feature\_group\_anfis\_clean.py}}\\[0.6em]
{\large Kiến trúc hai nhánh: ANFIS theo nhóm đặc trưng + BiLSTM theo chuỗi}
\end{center}

\begin{center}
\resizebox{0.98\linewidth}{!}{
\begin{tikzpicture}[node distance=8mm and 9mm]
  \node[box, fill=inputblue, minimum width=38mm] (input) at (0,0)
    {Đầu vào chuỗi\\\texttt{sequence\_input}\\$(B,\;look\_back,\;n\_seq\_features)$};

  \node[box, fill=fuzzygreen, minimum width=34mm] (lastcore) at (5.0,2.0)
    {Lấy bước cuối\\của 6 đặc trưng lõi\\\texttt{last\_core\_features}};
  \node[smallbox, fill=fuzzygreen, minimum width=31mm] (returns) at (9.4,3.1)
    {\texttt{returns\_slice}\\4 đặc trưng:\\\texttt{close\_ret}, \texttt{open\_gap},\\\texttt{high\_buffer}, \texttt{low\_buffer}};
  \node[smallbox, fill=fuzzygreen, minimum width=31mm] (indicators) at (9.4,0.9)
    {\texttt{indicator\_slice}\\2 đặc trưng:\\\texttt{range\_pct}, \texttt{volume\_ret}};
  \node[smallbox, fill=fuzzygreen, minimum width=27mm] (anfisr) at (13.5,3.1)
    {\texttt{ANFIS returns}\\\texttt{output\_dim = 4}\\số luật $= n\_mfs^4$};
  \node[smallbox, fill=fuzzygreen, minimum width=27mm] (anfisi) at (13.5,0.9)
    {\texttt{ANFIS indicators}\\\texttt{output\_dim = 4}\\số luật $= n\_mfs^2$};
  \node[box, fill=fuzzygreen, minimum width=27mm] (fconcat) at (17.4,2.0)
    {\texttt{fuzzy\_concat}};
  \node[smallbox, fill=fuzzygreen] (fhidden) at (20.8,2.0)
    {\texttt{Dense(16, tanh)}\\\texttt{fuzzy\_hidden}};
  \node[smallbox, fill=fuzzygreen] (fraw) at (24.1,2.0)
    {\texttt{Dense(4)}\\\texttt{fuzzy\_raw\_params}};

  \node[smallbox, fill=seqorange, minimum width=25mm] (bilstm1) at (5.2,-2.2)
    {\texttt{BiLSTM 1}\\\texttt{return\_sequences=True}};
  \node[smallbox, fill=seqorange] (drop1) at (8.5,-2.2)
    {\texttt{Dropout}};
  \node[smallbox, fill=seqorange] (bilstm2) at (11.5,-2.2)
    {\texttt{BiLSTM 2}};
  \node[smallbox, fill=seqorange] (drop2) at (14.4,-2.2)
    {\texttt{Dropout}};
  \node[smallbox, fill=seqorange] (shidden) at (17.6,-2.2)
    {\texttt{Dense(16, relu)}\\\texttt{sequence\_hidden}};
  \node[smallbox, fill=seqorange] (sraw) at (21.2,-2.2)
    {\texttt{Dense(4)}\\\texttt{sequence\_residual\_params}};

  \node[box, fill=mixpurple, minimum width=29mm] (add) at (27.8,0.0)
    {\texttt{raw\_target\_params}\\=\texttt{fuzzy\_raw} + \texttt{sequence\_raw}};

  \node[smallbox, fill=outred, minimum width=27mm] (close) at (32.6,2.6)
    {\texttt{close\_ret\_output}\\$RETURN\_CLIP \cdot \tanh(z_0)$};
  \node[smallbox, fill=outred, minimum width=27mm] (open) at (32.6,0.8)
    {\texttt{open\_gap\_output}\\$RETURN\_CLIP \cdot \tanh(z_1)$};
  \node[smallbox, fill=outred, minimum width=27mm] (high) at (32.6,-1.0)
    {\texttt{high\_buffer\_output}\\$BUFFER\_CLIP \cdot \sigma(z_2)$};
  \node[smallbox, fill=outred, minimum width=27mm] (low) at (32.6,-2.8)
    {\texttt{low\_buffer\_output}\\$BUFFER\_CLIP \cdot \sigma(z_3)$};
  \node[box, fill=outred, minimum width=31mm] (target) at (37.0,0.0)
    {\texttt{target\_output}\\$(close\_ret,\;open\_gap,$\\$high\_buffer,\;low\_buffer)$};
  \node[box, fill=outred, minimum width=33mm] (recon) at (41.9,0.0)
    {Dựng lại OHLC hợp lệ\\\texttt{reconstruct\_ohlc(...)}};

  \draw[line] (input) -- (lastcore);
  \draw[line] (lastcore) -- (returns);
  \draw[line] (lastcore) -- (indicators);
  \draw[line] (returns) -- (anfisr);
  \draw[line] (indicators) -- (anfisi);
  \draw[line] (anfisr) -- (fconcat);
  \draw[line] (anfisi) -- (fconcat);
  \draw[line] (fconcat) -- (fhidden);
  \draw[line] (fhidden) -- (fraw);

  \draw[line] (input) -- (bilstm1);
  \draw[line] (bilstm1) -- (drop1);
  \draw[line] (drop1) -- (bilstm2);
  \draw[line] (bilstm2) -- (drop2);
  \draw[line] (drop2) -- (shidden);
  \draw[line] (shidden) -- (sraw);

  \draw[line] (fraw) -- (add);
  \draw[line] (sraw) -- (add);
  \draw[line] (add) -- (close);
  \draw[line] (add) -- (open);
  \draw[line] (add) -- (high);
  \draw[line] (add) -- (low);
  \draw[line] (close) -- (target);
  \draw[line] (open) -- (target);
  \draw[line] (high) -- (target);
  \draw[line] (low) -- (target);
  \draw[line] (target) -- (recon);

  \node[note] at (20.7,-5.5) {
    \textbf{Ý tưởng chính:}\\
    Nhánh mờ học một tiên nghiệm có cấu trúc từ 6 đặc trưng lõi dễ diễn giải.\\
    Nhánh BiLSTM học phần hiệu chỉnh theo chuỗi thời gian còn thiếu.\\
    Hai nhánh được cộng ở mức \texttt{raw\_target\_params}, rồi 4 đầu ra cuối bị chặn về miền hợp lệ trước khi dựng lại OHLC.
  };
\end{tikzpicture}
}
\end{center}

\newpage

\begin{center}
{\LARGE\bfseries\color{titleblue} Sơ đồ chi tiết từ dữ liệu đến dự báo}\\[0.6em]
{\large Bản chi tiết bám sát các hàm \texttt{engineer\_features}, \texttt{build\_windows}, \texttt{build\_model}, \texttt{reconstruct\_ohlc}}
\end{center}

\begin{center}
\resizebox{0.99\linewidth}{!}{
\begin{tikzpicture}[node distance=7mm and 8mm]
  % preprocessing
  \node[box, fill=inputblue, minimum width=31mm] (csv) at (0,0) {CSV trong\\\texttt{Dataset/}};
  \node[smallbox, fill=inputblue, minimum width=34mm] (load) at (4.3,0)
    {\texttt{load\_market\_dataframe}\\- kiểm tra OHLC\\- parse \texttt{Date}\\- sort thời gian};
  \node[smallbox, fill=inputblue, minimum width=36mm] (feat) at (9.6,0)
    {\texttt{engineer\_features}\\- 6 đặc trưng lõi\\- 4 mục tiêu tham số hóa\\- lưu \texttt{next\_OHLC} thật};
  \node[smallbox, fill=inputblue, minimum width=39mm] (windows) at (15.6,0)
    {\texttt{build\_windows}\\- rolling windows\\- chia train/val/test theo thời gian\\- fit scaler chỉ trên train};

  % tensor summary
  \node[box, fill=inputblue, minimum width=39mm] (tensor) at (21.7,0)
    {Tensor đầu vào mô hình\\$X \in \mathbb{R}^{B \times T \times F}$\\và mục tiêu $y \in \mathbb{R}^{B \times 4}$};

  % fuzzy branch
  \node[groupbox, fit={(24.4,3.7) (38.2,9.8)}, label={[font=\bfseries\small]north west:Nhánh mờ}] (fuzzygroup) {};
  \node[smallbox, fill=fuzzygreen, minimum width=29mm] (last) at (26.6,8.3)
    {\texttt{last\_core\_features}\\$x[:, -1, :6]$};
  \node[smallbox, fill=fuzzygreen, minimum width=27mm] (rslice) at (31.1,9.0)
    {\texttt{returns\_slice}\\4 đặc trưng};
  \node[smallbox, fill=fuzzygreen, minimum width=27mm] (islice) at (31.1,7.0)
    {\texttt{indicator\_slice}\\2 đặc trưng};
  \node[smallbox, fill=fuzzygreen, minimum width=27mm] (ranfis) at (35.6,9.0)
    {\texttt{anfis\_returns}\\ra 4 chiều ẩn};
  \node[smallbox, fill=fuzzygreen, minimum width=27mm] (ianfis) at (35.6,7.0)
    {\texttt{anfis\_indicators}\\ra 4 chiều ẩn};
  \node[smallbox, fill=fuzzygreen, minimum width=26mm] (fcat) at (40.0,8.0)
    {\texttt{fuzzy\_concat}};
  \node[smallbox, fill=fuzzygreen, minimum width=26mm] (fdense) at (44.0,8.0)
    {\texttt{Dense(16, tanh)}};
  \node[smallbox, fill=fuzzygreen, minimum width=28mm] (fout) at (48.2,8.0)
    {\texttt{fuzzy\_raw\_params}\\$(B,4)$};

  % sequence branch
  \node[groupbox, fit={(24.4,-0.8) (41.4,4.2)}, label={[font=\bfseries\small]north west:Nhánh chuỗi}] (seqgroup) {};
  \node[smallbox, fill=seqorange, minimum width=26mm] (seqin) at (26.5,1.8)
    {\texttt{sequence\_input}\\toàn bộ cửa sổ};
  \node[smallbox, fill=seqorange, minimum width=24mm] (b1) at (30.3,1.8)
    {\texttt{BiLSTM 1}};
  \node[smallbox, fill=seqorange, minimum width=20mm] (d1) at (33.4,1.8)
    {\texttt{Dropout}};
  \node[smallbox, fill=seqorange, minimum width=24mm] (b2) at (36.3,1.8)
    {\texttt{BiLSTM 2}};
  \node[smallbox, fill=seqorange, minimum width=20mm] (d2) at (39.2,1.8)
    {\texttt{Dropout}};
  \node[smallbox, fill=seqorange, minimum width=28mm] (sdense) at (43.0,1.8)
    {\texttt{Dense(16, relu)}};
  \node[smallbox, fill=seqorange, minimum width=30mm] (sout) at (47.5,1.8)
    {\texttt{sequence\_residual\_params}\\$(B,4)$};

  % merge and outputs
  \node[groupbox, fit={(50.0,-0.8) (67.0,9.6)}, label={[font=\bfseries\small]north west:Kết hợp và đầu ra}] (mergegroup) {};
  \node[smallbox, fill=mixpurple, minimum width=30mm] (raw) at (53.1,4.8)
    {\texttt{raw\_target\_params}\\$=$ nhánh mờ + nhánh chuỗi};
  \node[smallbox, fill=outred, minimum width=24mm] (c1) at (57.6,7.7)
    {\texttt{close\_ret}\\$\tanh$};
  \node[smallbox, fill=outred, minimum width=24mm] (c2) at (57.6,5.8)
    {\texttt{open\_gap}\\$\tanh$};
  \node[smallbox, fill=outred, minimum width=24mm] (c3) at (57.6,3.9)
    {\texttt{high\_buffer}\\$\sigma$};
  \node[smallbox, fill=outred, minimum width=24mm] (c4) at (57.6,2.0)
    {\texttt{low\_buffer}\\$\sigma$};
  \node[smallbox, fill=outred, minimum width=30mm] (tout) at (62.2,4.8)
    {\texttt{target\_output}\\$(B,4)$};
  \node[smallbox, fill=outred, minimum width=33mm] (rohlc) at (66.2,4.8)
    {\texttt{reconstruct\_ohlc}\\dựng giá OHLC hợp lệ};

  % formulas
  \node[note] (formula) at (57.5,-3.2) {
    \textbf{Công thức dựng lại giá:}\\
    $pred\_close = current\_close \cdot (1 + close\_ret)$\\
    $pred\_open = current\_close \cdot (1 + open\_gap)$\\
    $body\_high = \max(pred\_open, pred\_close)$\\
    $body\_low = \min(pred\_open, pred\_close)$\\
    $pred\_high = body\_high \cdot e^{high\_buffer}$\\
    $pred\_low = body\_low \cdot e^{-low\_buffer}$
  };

  \node[note, text width=7.9cm] (guarantee) at (67.6,-3.2) {
    \textbf{Hệ quả cấu trúc:}\\
    Mô hình mới không dự đoán trực tiếp 4 giá độc lập.\\
    Vì vậy luôn có:\\
    $pred\_high \ge \max(pred\_open, pred\_close)$\\
    $pred\_low \le \min(pred\_open, pred\_close)$
  };

  % decompose model note
  \node[note, text width=11.0cm] (decomp) at (31.0,-3.2) {
    \textbf{Khối giải thích nội bộ \texttt{decompose\_model}:}\\
    Khi cần phân tích một mẫu dữ liệu, mô hình phụ này phơi bày đồng thời:\\
    1. \texttt{fuzzy\_raw\_params}\\
    2. \texttt{sequence\_residual\_params}\\
    3. \texttt{target\_output}\\
    Nhờ vậy ta nhìn thấy rõ nhánh mờ đóng góp gì và nhánh chuỗi đang sửa phần nào.
  };

  % arrows
  \draw[line] (csv) -- (load);
  \draw[line] (load) -- (feat);
  \draw[line] (feat) -- (windows);
  \draw[line] (windows) -- (tensor);

  \draw[line] (tensor) -- (last);
  \draw[line] (tensor) -- (seqin);

  \draw[line] (last) -- (rslice);
  \draw[line] (last) -- (islice);
  \draw[line] (rslice) -- (ranfis);
  \draw[line] (islice) -- (ianfis);
  \draw[line] (ranfis) -- (fcat);
  \draw[line] (ianfis) -- (fcat);
  \draw[line] (fcat) -- (fdense);
  \draw[line] (fdense) -- (fout);

  \draw[line] (seqin) -- (b1);
  \draw[line] (b1) -- (d1);
  \draw[line] (d1) -- (b2);
  \draw[line] (b2) -- (d2);
  \draw[line] (d2) -- (sdense);
  \draw[line] (sdense) -- (sout);

  \draw[line] (fout) -- (raw);
  \draw[line] (sout) -- (raw);
  \draw[line] (raw) -- (c1);
  \draw[line] (raw) -- (c2);
  \draw[line] (raw) -- (c3);
  \draw[line] (raw) -- (c4);
  \draw[line] (c1) -- (tout);
  \draw[line] (c2) -- (tout);
  \draw[line] (c3) -- (tout);
  \draw[line] (c4) -- (tout);
  \draw[line] (tout) -- (rohlc);
\end{tikzpicture}
}
\end{center}

\newpage

\begin{center}
{\LARGE\bfseries\color{titleblue} Sơ đồ chi tiết bên trong một khối \texttt{OrderedFeatureGroupANFIS}}\\[0.6em]
{\large Cấu trúc mà hai khối \texttt{anfis\_returns} và \texttt{anfis\_indicators} cùng dùng}
\end{center}

\begin{center}
\resizebox{0.95\linewidth}{!}{
\begin{tikzpicture}[node distance=8mm and 9mm]
  \node[box, fill=inputblue, minimum width=35mm] (inputa) at (0,0)
    {Đầu vào cục bộ\\\texttt{inputs}\\$(B,\;n\_features)$};

  \node[smallbox, fill=fuzzygreen, minimum width=33mm] (centers) at (5.0,2.0)
    {Tâm có thứ tự\\\texttt{center\_base}\\$+$ \texttt{softplus(center\_delta\_raw)}};
  \node[smallbox, fill=fuzzygreen, minimum width=33mm] (widths) at (5.0,-2.0)
    {Độ rộng Gaussian dương\\\texttt{softplus(width\_raw)}};
  \node[box, fill=fuzzygreen, minimum width=34mm] (fuzz) at (10.0,0.0)
    {Mờ hóa Gaussian\\\texttt{memberships[b, i, j]}};

  \node[smallbox, fill=fuzzygreen, minimum width=31mm] (indices) at (15.0,2.2)
    {\texttt{rule\_mf\_indices}\\luật nào dùng\\hàm thuộc nào};
  \node[box, fill=fuzzygreen, minimum width=34mm] (firing) at (15.0,-0.4)
    {Cường độ kích hoạt luật\\\texttt{firing}};
  \node[smallbox, fill=fuzzygreen, minimum width=30mm] (norm) at (20.0,-0.4)
    {Chuẩn hóa\\\texttt{firing\_norm}};

  \node[smallbox, fill=mixpurple, minimum width=34mm] (cons) at (20.0,3.0)
    {Hệ quả Sugeno\\\texttt{rule\_output}\\$= bias + \sum coeff_i x_i$};
  \node[box, fill=mixpurple, minimum width=35mm] (weighted) at (25.1,1.3)
    {Cộng có trọng số\\$\sum firing\_norm \cdot rule\_output$};
  \node[box, fill=outred, minimum width=34mm] (out) at (30.0,1.3)
    {Đầu ra ẩn của khối ANFIS\\$(B,\;output\_dim)$};

  \node[note, text width=12.0cm] (logic) at (16.0,-4.6) {
    \textbf{Diễn giải toán học:}\\
    Mỗi đặc trưng được mờ hóa bằng nhiều hàm thuộc Gaussian.\\
    Mỗi luật chọn đúng một hàm thuộc trên mỗi đặc trưng, rồi lấy tích để tạo cường độ kích hoạt.\\
    Sau khi chuẩn hóa, mỗi luật sinh ra một hệ quả affine kiểu Sugeno.\\
    Đầu ra cuối cùng của khối là tổng có trọng số của tất cả các hệ quả luật.
  };

  \node[note, text width=10.8cm] (count) at (30.3,-4.6) {
    \textbf{Số luật:}\\
    Nếu một khối có \texttt{n\_features} đầu vào và mỗi đầu vào có \texttt{n\_mfs} hàm thuộc thì:\\[0.2em]
    \[
      n\_rules = n\_mfs^{\,n\_features}
    \]
    Ví dụ:\\
    - \texttt{anfis\_returns}: $n\_features = 4 \Rightarrow n\_rules = n\_mfs^4$\\
    - \texttt{anfis\_indicators}: $n\_features = 2 \Rightarrow n\_rules = n\_mfs^2$
  };

  \draw[line] (inputa) -- (fuzz);
  \draw[line] (centers) -- (fuzz);
  \draw[line] (widths) -- (fuzz);
  \draw[line] (fuzz) -- (firing);
  \draw[line] (indices) -- (firing);
  \draw[line] (firing) -- (norm);
  \draw[line] (inputa) |- (cons);
  \draw[line] (norm) -- (weighted);
  \draw[line] (cons) -- (weighted);
  \draw[line] (weighted) -- (out);
\end{tikzpicture}
}
\end{center}

\end{document}
